\section{Background}

The proliferation of autonomous multi-agent AI systems has exposed a fundamental gap between the speed of machine decision-making and the capacity of human oversight to keep pace. As agentic systems grow more capable of chaining actions across software interfaces, filesystems, and network endpoints, the question of who—or what—bears responsibility for downstream consequences has become both a technical and a regulatory imperative. This section surveys the prior art across accountability architectures, human-in-the-loop design, and fail-safe engineering, then positions the BX3 Framework's three-layer model as an integrative response to the failure modes that prior approaches have left unaddressed.

% -------------------------------------------------------------------
\subsection{Accountability in Multi-Agent AI Systems}

Classical accountability frameworks borrowed from organizational theory assume that agents are human actors with legible motivations and bounded liability \cite{wiener1948}. When an AI system acts autonomously, the causal chain between intention and outcome becomes opaque: a model may produce a recommendation that triggers a cascade of downstream system behaviors, none of which were explicitly programmed or foreseen by the model's operators. This opacity is compounded in multi-agent settings, where one agent's output becomes another agent's input, creating distributed chains of deferred responsibility \cite{bansal2019}.

Dijkstra's seminal work on cooperative sequential processes offered an early formal lens for understanding how discrete, independent actors can be composed into larger computing systems without sacrificing logical rigor \cite{dijkstra1982}. While his concern was with correctness proofs rather than accountability, the insight that composability demands explicit interfaces translates directly to modern agentic pipelines: each agent-to-agent handoff requires a formally specifiable contract describing preconditions, postconditions, and failure modes. Without such contracts, accountability dissolves into the gap between actors.

The AI governance literature has proposed various accountability mechanisms—audit trails, model cards, and explainability reports—but these operate primarily at the model level, not at the workflow level. They assume a single deployer-single model relationship, which breaks down in multi-agent orchestration where many specialized models interact dynamically. A robust accountability model must therefore span the entire execution graph, recording not just model outputs but also which human principals authorized which agent actions and at what escalation threshold \cite{tristr81}.

% -------------------------------------------------------------------
\subsection{Human-in-the-Loop Design}

The human-in-the-loop (HITL) paradigm originated in safety-critical control systems, where the operator remains the ultimate authority over any automated decision that could result in harm. Classical HITL design treats automation as a decision support tool: the machine proposes, the human approves or overrides. This model works well in low-frequency, high-stakes contexts (e.g., nuclear plant control rooms) but degrades under high-frequency, high-volume conditions where the pace of machine events exceeds human cognitive bandwidth.

Prior research distinguishes three HITL modes: full human control (manual), human-on-the-loop (monitoring with override authority), and human-off-the-loop (fully autonomous). Most AI agent frameworks today operate in an implicit human-off-the-loop mode for routine sub-tasks while retaining ad-hoc human escalation for exceptional cases. This hybrid is insufficient for agentic systems precisely because the boundary between routine and exceptional is not known \textit{a priori} \cite{bansal2019}. An agent that decides to delete a file, send an email, or charge a credit card may be operating within its programmed scope, yet the downstream consequences of such actions can cross safety thresholds without triggering any human review.

Effective HITL design in agentic contexts must therefore be \textit{tiered by consequence severity}, not by action frequency. Actions that are irreversible, costly, or access private data should require explicit human confirmation before execution, regardless of how routine they appear to the agent's planning logic. This requires a formal classification schema that maps action types to required human involvement levels—a schema that the BX3 Framework operationalizes through its DAP layer.

% -------------------------------------------------------------------
\subsection{Fail-Safe Design in Safety-Critical Systems}

The concept of fail-safe design holds that systems should degrade in a manner that preserves human safety above all other concerns. In classical control theory, this is often achieved through watchdog timers, heartbeat signals, and authority limiting—mechanisms that ensure no single automated component can place the system in an irrecoverable state without human intervention \cite{wiener1948}. The fail-safe principle is particularly well-established in aviation and medical device software, where regulatory standards such as DO-178C and IEC 62304 mandate specific error handling and recovery procedures.

Two failure patterns are especially relevant to agentic AI systems. The first is \textbf{cascade failure}, in which a fault in one agent propagates to downstream agents that trust its output, eventually contaminating large portions of the execution graph with a single erroneous decision \cite{tristr81}. The second is \textbf{silent failure}, in which an agent continues to execute without signaling an anomalous condition because its error-handling logic was never exercised during training or testing \cite{bansal2019}. Both patterns are exacerbated by the absence of formal correctness proofs at agent handoff boundaries.

Classical fail-safe approaches rely on containment strategies: isolating faulty subsystems, reverting to known-good states, and providing human operators with sufficient information to assess the situation and take manual control \cite{dijkstra1982}. In software systems, this is implemented through transaction rollback, circuit breakers, and watchdog processes. Translating these patterns to agentic workflows requires an analogous recovery architecture—one that can identify which agent actions are irrecoverable, preserve sufficient state to reason about the fault, and provide the human operator with a meaningful override capability before harm propagates \cite{wiener1948}.

The challenge unique to AI agents is that the failure modes are partially stochastic and context-dependent. A model that fails to recognize a safety-critical condition in one context may succeed in another, making it impossible to enumerate all failure modes at design time. This inherent unpredictability demands a runtime monitoring and bailout architecture that does not depend on exhaustive pre-specification of failure modes \cite{tristr81}.

% -------------------------------------------------------------------
\subsection{The BX3 Framework Three-Layer Model}

The BX3 Framework \cite{bansal2019} addresses the aforementioned gaps through a three-layer architectural model that separates concerns across the Execution Plane, the Orchestration Plane, and the Audit Plane. Each layer has a distinct responsibility and a formally defined interface to the others, enabling composability, inspectability, and fail-safe recovery.

The \textbf{Execution Plane (XP)} is the lowest layer and is responsible for the actual execution of agent tasks. Each agent in the XP operates within a bounded, independently verifiable execution context. The XP exposes a standardized failure-signal interface that propagates error conditions upward to the Orchestration Plane. Critically, the XP enforces a non-negotiable constraint: no execution step may permanently alter external state (file deletion, financial transaction, network transmission) without a positive acknowledgment from the layer above. This constraint operationalizes the fail-safe principle at the execution level.

The \textbf{Decision Authority Plane (DAP)} is the middle layer and serves as the locus of human oversight and authorization. The DAP maintains a structured policy schema that maps action types to required authorization levels. Actions classified as high-consequence require explicit human confirmation; low-consequence actions may proceed under automated authorization within defined parameter bounds. The DAP also manages escalation—receiving failure signals from the XP and routing them to the appropriate human or automated response handler. This layer is where the HITL paradigm is made concrete: it is not a passive observer but an active gatekeeper that can halt, redirect, or abort execution flows before harm propagates.

The \textbf{Audit Plane (AP)} is the top layer and serves as the forensic and accountability layer. The AP records every execution event, every authorization decision, every human override, and every failure signal in a tamper-evident log. The AP does not intervene in execution but provides the data necessary for post-hoc accountability: determining which policy was in effect, what the human operator authorized, what the agent actually executed, and where the causal chain diverged from expectation. This log serves both operational debugging and external regulatory compliance.

The BX3 model is distinguished from prior architectural approaches in its explicit treatment of the human as a first-class component within the control loop, not as an afterthought or a fallback \cite{tristr81}. By separating execution (XP), authorization (DAP), and accountability (AP) into distinct planes with formal interfaces, the framework enables each concern to be reasoned about, verified, and updated independently—without requiring the entire system to be re-certified when a policy changes or an agent is updated.

The three-layer model also provides a natural locus for implementing the bailout operation. When a fail-safe condition is detected—either by runtime monitoring or by human observation—the bailout protocol operates through the DAP to halt the XP agents, preserve execution state for forensic analysis, and return the system to a known-safe configuration. This protocol is the subject of the present paper.

% -------------------------------------------------------------------
\subsection{Summary}

The foregoing review establishes that accountability, HITL oversight, and fail-safe design are not independent concerns but interrelated requirements for trustworthy agentic systems. Prior work has addressed each in isolation, but the integration of all three into a coherent architectural model remains an open problem. The BX3 Framework's three-layer model offers a principled structure for addressing this integration. The following section formalizes the Bailout Protocol as a specific mechanism operating within the DAP layer of that model.